\section{Related Work}
\label{sec:related_work}

Our framework interfaces with three paradigms in software automation: automated program repair, code-centric retrieval, and collaborative multi-agent system design.

\subsection{Automated Program Repair}

Traditional automated program repair (APR) systems rely on search-based heuristics, genetic programming, templates, or semantic constraints to manipulate programs under a test or specification signal. These methods can provide strong precision within localized settings, but often depend on concrete test suites or other executable specifications~\cite{legoues2019apr}. SWE-bench reframed this challenge for language-model agents by evaluating whether a model can edit real repositories in response to natural-language issues~\cite{jimenez2024swebench}. The Context Sphere builds on this problem setting, but focuses on the context-orchestration layer: how an agent should select, expand, and reuse repository evidence before and during patch generation.

\subsection{Repository Retrieval Pipelines}

Codebase retrieval systems commonly map natural-language issues to files, functions, or chunks using lexical retrieval, embeddings, reranking, or learned selectors. This makes large repositories tractable, but standard retrieval can ignore the dependency topology that makes a code edit valid. RepoCoder demonstrates an iterative retrieval-generation loop for repository-level code completion, while long-context studies show that adding more context is not always sufficient because models can fail to use information buried in the middle of a prompt~\cite{zhang2023repocoder,liu2024lost}. Graph-structured code models and graph-based repository retrieval methods further motivate topology-aware context, but they often require richer graph construction or graph-aware representation learning~\cite{guo2021graphcodebert,liu2024graphcoder}. The Context Sphere occupies a narrower and more lightweight point in this design space: a neural selector identifies the likely issue centroid, and a deterministic AST expansion step restores immediate local dependencies.

\subsection{Multi-Persona LLM Architectures}

Multi-agent software systems such as ChatDev, MetaGPT, and SWE-agent demonstrate that LLMs can be organized into role-specific workflows for design, coding, testing, repository navigation, and tool use~\cite{qian2024chatdev,hong2024metagpt,yang2024sweagent}. These systems motivate role separation, but open-ended dialogue can accumulate prompt noise and confirmation bias. Sycophancy studies show that assistant models can prefer responses that align with user beliefs over truthful correction, motivating stricter reviewer boundaries~\cite{sharma2023sycophancy}. The Context Sphere differs by coupling role separation to a bounded repository neighborhood and by forcing Worker--Reviewer iteration through a structured constraint ledger.

\subsection{Cost-Aware Agent Deployment}

End-to-end coding agents are increasingly evaluated as systems rather than isolated generation calls. This motivates cost-aware reporting alongside correctness, since long-horizon repository repair can require many retrieval, editing, review, and verification turns. HAL and SWE-Effi both reflect this shift by surfacing resource consumption together with task performance~\cite{princetonhal2026,fan2025sweeffi}. Our work follows this direction by reporting local repair success together with token and cost estimates, and by treating context selection as an optimization target rather than an implementation detail.

\subsection{LLM-as-Judge Reliability}

Although the primary ablation reports execution-checked outcomes under a local harness, LLM-as-judge reliability remains relevant because the orchestration loop still uses a Reviewer persona before verification. Prior work on LLM-based judging has documented position, verbosity, and self-enhancement biases~\cite{zheng2023judging}. Accordingly, this paper treats internal approval as an intermediate control signal. Official benchmark execution, gold-patch comparison, and swapped-order judging remain necessary before claiming leaderboard-level repair correctness.
